\section{Processing the period groups}\label{sec:processing}

We now give the connection procedure. Process the primitive-text groups in increasing order of least period, breaking ties by the lexicographically least rotation of their primitive texts. Requests reserve budgets and are fulfilled when their host group is processed. Ordinary links are planned now and finalized in \Cref{sec:cycles}. Every record and every layer in these plans refers to the fixed base decomposition from \Cref{lem:layers}. For bookkeeping, mark a layer rooted when a stated operation explicitly roots it. Incidental connections to $\eps$ do not change these marks or budget assignments.

\subsection{Incoming requests and a baseline}

At entry to a group of period $p$, fulfill all requests received by each recipient layer in one application of \Cref{lem:host}. We will prove below that every requester has smaller period, and that a child group sends requests using only one record on one recipient layer. These are precisely the facts needed to apply that lemma with fresh recipient budgets.

Every recipient has an actual vertex of length less than $2p$: its record with a child of period $r<p$ has length less than $p+r$ by \Cref{lem:agreement}. If there are recipients, choose their highest layer as a \emph{baseline}. They and their requesters have already been rooted.

Each remaining layer below the baseline can be rooted with its own budget $p$. If it has a vertex of length at most $p$, take a round trip through $\eps$. Otherwise choose a baseline vertex $S=[a,h)$ of length less than $2p$. A lower layer has a vertex $[a,e)$ with
\[
p<e-a\le h-a<2p.
\]
For example, take $e=\min(l_i(a),h)$; pointwise ordering ensures this lies in its vertical range. The windows
\[
[a,h),\qquad[a+p,e+p),\qquad[a+p,h+p)
\]
are ordered because $h<a+2p<e+p$. Their cost-$p$ closed walk touches the lower layer and the rooted baseline. Recipient layers themselves need no further budget.

Let $n$ be the number of layers above the baseline, or the entire number in the group when there is no baseline. Their budgets are all still available. If $n=0$, proceed to the next group.

\subsection{Easy connections for the upper layers}

Use a distinguished position $t$ of the group's text $A$, and put
\[
H=f_H(t),
\]
the top layer's first end at that start. If $H-t\le np$, a round trip from $\eps$ through $[t,H)$ costs at most $np$. Its up path touches every upper layer, since every such layer has first end at most $H$. This roots them all.

If a baseline exists and its last end $E$ at $t$ satisfies $E\ge H-np$, \Cref{lem:band}, including the baseline as the bottom layer and taking $h=np$, joins it to all upper layers within their combined budget. This also roots them all.

It remains to handle the case
\begin{equation}\label{eq:hard-case}
H>t+np,
\qquad l_{\mathrm{base}}(t)<H-np
\quad\text{when a baseline exists}.
\end{equation}

\subsection{Forcing a record on another aligned text}

The count rule will supply $k+1$ unblocked occurrences of a word obtained by shortening the top vertex by $kp$. We choose $k$ so that at most $k$ of those occurrences can have the present full-text alignment. The same choice must keep the top $k$ layers close enough to join for at most $(k-1)p$, leaving one period of their combined budget available. The following threshold search gives both properties.

Starting at $k=1$, increase $k$ while the number of layers of this group satisfying
\[
l_i(t)\ge H-kp
\]
exceeds $k$. Stop at the first failure of this test and set $K=kp$.

\begin{lemma}\label{lem:choose-k}
This procedure stops with $1\le k\le n$. If $L_*$ is the lowest of the top $k$ layers, then
\begin{align}
l_{L_*}(t)&\ge H-(k-1)p,\label{eq:sharp-band}\\
f_H(x)-l_{L_*}(x)&\le K\qquad(x\in\ZZ).\label{eq:wide-band}
\end{align}
At most $k$ recorded occurrences of $w=A[t:H-K)$ have a full-text alignment identical to $A$.
\end{lemma}

\begin{proof}
For $k\le n$, layers at or below the baseline do not qualify, by \eqref{eq:hard-case}. Thus the process stops by $n$. If $k=1$, the top last end is at least its first end $H$. If $k>1$, reaching $k$ means that at least $k$ layers qualified at the previous threshold $H-(k-1)p$. Pointwise ordering gives \eqref{eq:sharp-band}.

For any $x$, choose its translate $x'$ modulo $p$ in $(t-p,t]$. Then
\[
f_H(x')\le H,
\qquad l_{L_*}(x')\ge l_{L_*}(t)-p.
\]
Their difference is invariant under period translation, proving \eqref{eq:wide-band}.

By \eqref{eq:hard-case}, $H-K>t$, so $w$ is nonempty. Each primitive layer has exactly one position residue whose full-text alignment agrees with $A$: the shifts preserving its text are exactly the multiples of $p$. At that alignment, $w$ is recorded exactly when $l_i(t)\ge H-K$. At most $k$ layers satisfy this inequality when the test stops. Other primitive-text groups have no identical full-text alignment.
\end{proof}

Apply the count rule \eqref{eq:period-rule}, with origin shifted to $t$, to the actual vertex word
\[
v=A[t:H),\qquad w=A[t:H-K).
\]
The top layer records an unblocked occurrence of $v$ in its own text. By \eqref{eq:unblocked-count}, the trigger holds, and the rule gives at least $k+1$ unblocked occurrences of $w$. \Cref{lem:choose-k} leaves at least one on a layer $D$ whose aligned text differs from $A$. A different phase of this same group is allowed.

Choose one such occurrence. \Cref{lem:matching-window} gives an actual window $R=[a,b)$ on $D$, agreeing with $A$ throughout and containing the specified interval $[t,H-K)$. Let $q$ be the least period of $D$. Thus
\begin{equation}\label{eq:chosen-record}
a\le t,
\qquad b\ge H-K,
\qquad b-a<p+q.
\end{equation}
If $D$ belongs to this same primitive-text group, then
\begin{equation}\label{eq:same-group-record}
b-a<p.
\end{equation}
The strict length bounds follow from \Cref{lem:agreement}. Use this \emph{one} chosen record for every subsequent connection from this processed group.

\subsection{Distributing the remaining budgets}

At start $a$, put $r_H=f_H(a)$ and $r_L=l_{L_*}(a)$. Monotonicity, \eqref{eq:wide-band}, and \eqref{eq:chosen-record} give
\begin{equation}\label{eq:block-preconditions}
r_H-r_L\le K,
\qquad r_H\le H\le b+K.
\end{equation}

\paragraph{The collective case.}
If $r_L+K\ge b$, make the top $k$ layers one \emph{collective block}. Plan the link in \Cref{lem:collective} from these layers to $D$ at $R$, reserving their combined budget $K$. Here a collective block is allowed to have $k=1$; its plan is still the collective construction. If the test fails, do not make the block. In that event,
\begin{equation}\label{eq:failed-collective}
r_H\le r_L+K<b.
\end{equation}

\paragraph{The individual cases.}
Handle every upper layer not placed in that collective block separately. Its first end at $a$ is at most $b$. When the collective test failed, this follows from \eqref{eq:failed-collective}; a layer below the top $k$ has
\[
f_i(a)\le l_i(t)<H-K\le b
\]
by the stopping rule and $a\le t$.

For such a layer, do exactly one of the following.
\begin{enumerate}[label=\textup{(\roman*)}]
\item If some vertex has length at most $p$, root the layer by a round trip through $\eps$ using at most its budget.
\item If all its vertex lengths exceed $p$ and
\begin{equation}\label{eq:single-test}
l_i(a)+p\ge b,
\end{equation}
make it an \emph{individual block} with a planned link to $D$ at $R$. At start $a+p$, its vertical range intersects $[b,b+p]$, since
\[
f_i(a+p)=f_i(a)+p\le b+p,
\qquad l_i(a+p)=l_i(a)+p\ge b.
\]
Choose an actual end $e'\in[b,b+p]$ there and follow
\begin{equation}\label{eq:individual-list}
[a,b),\qquad[a+p,e'),\qquad[a+p,b+p).
\end{equation}
This is a cost-$p$ closed walk in $A$ touching the layer. Its incoming transition to start $a+p$ reaches a word of length at most
\begin{equation}\label{eq:single-short}
\max(0,b-a-p)\le q,
\end{equation}
by \eqref{eq:chosen-record} and \Cref{lem:ordered-windows}. Reserve this layer's budget for that link.
\item Otherwise reserve the layer's budget $p$ and send a request to $D$ at the same record $R$. Its first end $e=f_i(a)$ satisfies
\begin{equation}\label{eq:outgoing-request}
a+p<e\le l_i(a)<b-p.
\end{equation}
This is the condition of \Cref{lem:host}. In particular $b-a>2p$; combining this with $b-a<p+q$ gives $q>p$. The recipient group will therefore be processed later.
\end{enumerate}

This verifies the promised hypotheses for incoming requests. They always go to strictly larger periods, and a child group uses just one record on one recipient layer, even if several of its layers request there. Ordinary planned links consume no recipient budget. When a recipient is processed, its own budget and all reserved requester budgets are available. After \Cref{lem:host} is applied, those layers are rooted and do not enter any outgoing block.

\begin{proposition}\label{prop:partition}
After all groups have been processed, every layer is either marked rooted or belongs to exactly one outgoing block. These two classes are disjoint. Each block consists of layers from one primitive-text group and targets a single actual layer. Its total budget remains reserved for its planned link. A collective block of $k$ layers and period $p$ also satisfies \eqref{eq:sharp-band} at its distinguished position; an individual block consists of one layer.
\end{proposition}

\begin{proof}
Incoming requests use disjoint requester budgets, and each recipient budget is spent only when its group is processed. The baseline step uses only budgets of remaining lower layers. The easy upper-layer branches use at most the combined budgets of the upper layers and finish them. In the remaining branch, the collective case and the individual cases partition those upper layers. Each request is eventually fulfilled, because its recipient is processed later. The layers not marked rooted by a direct operation or a fulfilled request are precisely the outgoing blocks. Their links have only been planned, so their budgets are unspent. The asserted inequality for collective blocks is \eqref{eq:sharp-band}.
\end{proof}
